% GENERATED FILE. Edit canonical lessons, then run npm run generate:books.
% canonical-source-hash: fnv1a64:734d51cd120dca36
% canonical-lessons: JA-C303-ryuu, JA-C303-numa, JA-C303-ukabu, JA-C303-hokora, JA-C303-medetai

\chapter{The Dragon in the Marsh}
\label{ch:ja-the-dragon-in-the-marsh}
\hlchaptermodality{303}

\begin{chapteropening}
\textbf{By the end of this chapter:} \emph{I can say where the dragon lived and what floated there.}

Complete the last lesson of chapter 303: I can say where the dragon lived and what floated there.
\end{chapteropening}

\section[\texorpdfstring{\ja{りゅう} (ryū)}{りゅう (ryū)}]{\ja{りゅう} (ryū) — a dragon}

\label{lesson:JA-C303-ryuu}



\begin{quote}
Before the new one: say the Japanese for a legend, then the Japanese for a god.
\end{quote}

\subsection*{You'll want to know: \ja{りゅう}}
\textbf{\ja{りゅう}} — \emph{ryū} — \textquotedblleft{}a dragon\textquotedblright{}.

\begin{quote}
\textbf{\ja{むかし、りゅうが} \ja{そらを} \ja{とんでいました。}} — \emph{mukashi, ryū ga sora o tonde imashita} — \textquotedblleft{}Long ago a dragon flew in the sky.\textquotedblright{}
\end{quote}

\begin{cousinweb}
\textbf{\ja{りゅう}} (\emph{ryū}) is Sino-Japanese (kango): one Chinese character read with its on'yomi, a reading borrowed from Chinese. It is \emph{ryū}, \textquotedblleft{}dragon\textquotedblright{} (Middle Chinese \emph{ljowng}) (KANJIDIC2; Middle Chinese in Baxter's transcription).
\end{cousinweb}

\begin{grammarlens}[title={What you've built}]
One more word for telling what happened.
\end{grammarlens}

\subsection*{Your turn}
\begin{itemize}
\raggedright
  \item \emph{Say it:} \emph{ryū}
  \item \emph{Say it:} \emph{ryū}, once more
  \item \emph{Say it:} \emph{mukashi, ryū ga sora o tonde imashita} — Long ago a dragon flew in the sky.
  \item \emph{Recall:} say the Japanese for a legend, then the Japanese for a god, then say \emph{ryū} again
\end{itemize}

\begin{tcolorbox}[breakable,colback=teal!4,colframe=teal!35!black,title={Before you move on}]
What does \ja{りゅう} mean? (\textquotedblleft{}A dragon\textquotedblright{}.) Say it once more.
\end{tcolorbox}
\section[\texorpdfstring{\ja{ぬま} (numa)}{ぬま (numa)}]{\ja{ぬま} (numa) — a marsh; a pond}

\label{lesson:JA-C303-numa}



\begin{quote}
Before the new one: say the Japanese for a god, then the Japanese for a dragon.
\end{quote}

\subsection*{You'll want to know: \ja{ぬま}}
\textbf{\ja{ぬま}} — \emph{numa} — \textquotedblleft{}a marsh; a pond\textquotedblright{}.

\begin{quote}
\textbf{\ja{りゅうは} \ja{ぬまに} \ja{すんでいました。}} — \emph{ryū wa numa ni sunde imashita} — \textquotedblleft{}The dragon lived in a marsh.\textquotedblright{}
\end{quote}

\begin{cousinweb}
\textbf{\ja{ぬま}} (\emph{numa}) is a native Japanese word (wago), inherited rather than borrowed from Chinese; its kanji was attached later as a kun'yomi (JMdict).
\end{cousinweb}

\begin{grammarlens}[title={What you've built}]
One more word for telling what happened.
\end{grammarlens}

\subsection*{Your turn}
\begin{itemize}
\raggedright
  \item \emph{Say it:} \emph{numa}
  \item \emph{Say it:} \emph{numa}, once more
  \item \emph{Say it:} \emph{ryū wa numa ni sunde imashita} — The dragon lived in a marsh.
  \item \emph{Recall:} say the Japanese for a god, then the Japanese for a dragon, then say \emph{numa} again
\end{itemize}

\begin{tcolorbox}[breakable,colback=teal!4,colframe=teal!35!black,title={Before you move on}]
What does \ja{ぬま} mean? (\textquotedblleft{}A marsh; a pond\textquotedblright{}.) Say it once more.
\end{tcolorbox}
\section[\texorpdfstring{\ja{うかぶ} (ukabu)}{うかぶ (ukabu)}]{\ja{うかぶ} (ukabu) — to float up; to come to mind}

\label{lesson:JA-C303-ukabu}



\begin{quote}
Before the new one: say the Japanese for a dragon, then the Japanese for a marsh.
\end{quote}

\subsection*{You'll want to know: \ja{うかぶ}}
\textbf{\ja{うかぶ}} — \emph{ukabu} — \textquotedblleft{}to float up; to come to mind\textquotedblright{}.

\begin{quote}
\textbf{\ja{ぬまに} \ja{とりが} \ja{うかんでいます。}} — \emph{numa ni tori ga ukande imasu} — \textquotedblleft{}Birds are floating on the marsh.\textquotedblright{}
\end{quote}

Polite \textbf{\ja{うかびます}}, \emph{ukabimasu}; a godan verb. Past (polite): \textbf{\ja{ふねが} \ja{うかびました。}}, \emph{fune ga ukabimashita}.

\begin{cousinweb}
\textbf{\ja{うかぶ}} (\emph{ukabu}) is a native Japanese word (wago), inherited rather than borrowed from Chinese; its kanji was attached later as a kun'yomi (JMdict).
\end{cousinweb}

\begin{grammarlens}[title={What you've built}]
One more word for telling what happened.
\end{grammarlens}

\subsection*{Your turn}
\begin{itemize}
\raggedright
  \item \emph{Say it:} \emph{ukabu}
  \item \emph{Say it:} \emph{ukabu}, once more
  \item \emph{Say it:} \emph{numa ni tori ga ukande imasu} — Birds are floating on the marsh.
  \item \emph{Recall:} say the Japanese for a dragon, then the Japanese for a marsh, then say \emph{ukabu} again
\end{itemize}

\begin{tcolorbox}[breakable,colback=teal!4,colframe=teal!35!black,title={Before you move on}]
What does \ja{うかぶ} mean? (\textquotedblleft{}To float up; to come to mind\textquotedblright{}.) Say it once more.
\end{tcolorbox}
\section[\texorpdfstring{\ja{ほこら} (hokora)}{ほこら (hokora)}]{\ja{ほこら} (hokora) — a small wayside shrine}

\label{lesson:JA-C303-hokora}



\begin{quote}
Before the new one: say the Japanese for a marsh, then the Japanese for to float up.
\end{quote}

\subsection*{You'll want to know: \ja{ほこら}}
\textbf{\ja{ほこら}} — \emph{hokora} — \textquotedblleft{}a small wayside shrine\textquotedblright{}.

\begin{quote}
\textbf{\ja{ぬまの} \ja{ちかくに} \ja{ちいさい} \ja{ほこらが} \ja{あります。}} — \emph{numa no chikaku ni chiisai hokora ga arimasu} — \textquotedblleft{}There is a little shrine near the marsh.\textquotedblright{}
\end{quote}

\begin{cousinweb}
\textbf{\ja{ほこら}} (\emph{hokora}) is a native Japanese word (wago), inherited rather than borrowed from Chinese; its kanji was attached later as a kun'yomi (JMdict).
\end{cousinweb}

\begin{grammarlens}[title={What you've built}]
One more word for telling what happened.
\end{grammarlens}

\subsection*{Your turn}
\begin{itemize}
\raggedright
  \item \emph{Say it:} \emph{hokora}
  \item \emph{Say it:} \emph{hokora}, once more
  \item \emph{Say it:} \emph{numa no chikaku ni chiisai hokora ga arimasu} — There is a little shrine near the marsh.
  \item \emph{Recall:} say the Japanese for a marsh, then the Japanese for to float up, then say \emph{hokora} again
\end{itemize}

\begin{tcolorbox}[breakable,colback=teal!4,colframe=teal!35!black,title={Before you move on}]
What does \ja{ほこら} mean? (\textquotedblleft{}A small wayside shrine\textquotedblright{}.) Say it once more.
\end{tcolorbox}
\section[\texorpdfstring{\ja{めでたい} (medetai)}{めでたい (medetai)}]{\ja{めでたい} (medetai) — auspicious; joyful}

\label{lesson:JA-C303-medetai}



\begin{quote}
Before the new one: say the Japanese for to float up, then the Japanese for a small wayside shrine.
\end{quote}

\subsection*{You'll want to know: \ja{めでたい}}
\textbf{\ja{めでたい}} — \emph{medetai} — \textquotedblleft{}auspicious; joyful\textquotedblright{}.

\begin{quote}
\textbf{\ja{あにの} \ja{けっこんは} \ja{めでたいです。}} — \emph{ani no kekkon wa medetai desu} — \textquotedblleft{}My brother's wedding is a happy event.\textquotedblright{}
\end{quote}

An i-adjective.

\begin{cousinweb}
\textbf{\ja{めでたい}} (\emph{medetai}) is a native Japanese word (wago), inherited rather than borrowed from Chinese, and it is usually written in kana (JMdict).
\end{cousinweb}

\begin{grammarlens}[title={What you've built}]
One more word for telling what happened.
\end{grammarlens}

\subsection*{Your turn}
\begin{itemize}
\raggedright
  \item \emph{Say it:} \emph{medetai}
  \item \emph{Say it:} \emph{medetai}, once more
  \item \emph{Say it:} \emph{ani no kekkon wa medetai desu} — My brother's wedding is a happy event.
  \item \emph{Recall:} say the Japanese for to float up, then the Japanese for a small wayside shrine, then say \emph{medetai} again
\end{itemize}

\begin{tcolorbox}[breakable,colback=teal!4,colframe=teal!35!black,title={Before you move on}]
What does \ja{めでたい} mean? (\textquotedblleft{}Auspicious; joyful\textquotedblright{}.) Say it once more.
\end{tcolorbox}
